\scanpage{71}
Ces suites exactes déterminent $G,\widetilde G,\mathfrak g,\widetilde{\mathfrak g}$ chacun directement en termes de $\mathcal L(M)$ et $\det$ (pour les groupes) et $M$ et $\operatorname{tr}$ (pour les algèbres de Lie). Bien entendu, $M$ est considéré comme algèbre de Lie \unsure{grâce} à la forme \unsure{trace}.

Les hom canoniques
\[
\left\{\begin{aligned}
\widetilde G&\longrightarrow G\\
\widetilde{\mathfrak g}&\xrightarrow{f}\mathfrak g
\end{aligned}\right.\tag*{(36}
\]
sont les composés
\[
\widetilde G\overset{\mathrm{incl.}}{\lhook\joinrel\longrightarrow}\mathcal L(M)\xrightarrow{\mathrm{proj}}G\tag*{(37)}
\]
\[
\widetilde{\mathfrak g}\overset{\mathrm{incl}}{\lhook\joinrel\longrightarrow}M\xrightarrow{\mathrm{proj}}\mathfrak g\tag*{(38)}
\]

Considérons sur $M$ la forme \deleted{quadratique} \inserted{bilinéaire} \deleted{\unsure{obtenue}}.
\[
\deleted{\unsure{\Phi|q}}\Phi(uv)=\operatorname{Tr}(uv)\tag*{(39)}
\]
On voit alors que $\varphi$ est symétrique et \unsure{$f$|$\varphi$} non dégénérée en chaque fibre, permettant donc d'identifier $M$ à son propre dual
\[
M\simeq M^\vee\tag*{(40)}
\]
L'orthogonale de $\mathcal O_S$ est $\widetilde{\mathfrak g}$, et on trouve ainsi un accouplement dualisant
% archive sheet number: 70
